\def\tituloconcepto{Algoritmo de Grover}
% Si el capitulo se compila solo (sin main.tex), esta cabecera abre el documento.
% Antes de % Si el capitulo se compila solo (sin main.tex), esta cabecera abre el documento.
% Antes de % Si el capitulo se compila solo (sin main.tex), esta cabecera abre el documento.
% Antes de \input{cabecera} cada capitulo define \tituloconcepto.
\ifdefined\enmain\else
  \documentclass[11pt,a4paper]{article}
  \input{preambulo}
  \begin{document}
  \title{\tituloconcepto}
  \author{\autor}
  \date{\today}
  \maketitle
\fi
 cada capitulo define \tituloconcepto.
\ifdefined\enmain\else
  \documentclass[11pt,a4paper]{article}
  % Preambulo compartido por todos los documentos. Compilar con XeLaTeX/LuaLaTeX
% (o Tectonic). Si se usa Overleaf, elegir el compilador XeLaTeX.
\usepackage[spanish,es-noshorthands]{babel}
\usepackage{amsmath,amssymb}
\usepackage{booktabs,array}
\usepackage[a4paper,margin=2.5cm]{geometry}
\usepackage{xcolor}
\usepackage{listings}
\usepackage[most]{tcolorbox}
\usepackage{quantikz}
\usepackage[hidelinks]{hyperref}

% ---- Autor (editar a gusto) -------------------------------------------------
\newcommand{\autor}{Genesis Silesky}

% ---- Notacion de Dirac y atajos ---------------------------------------------
% quantikz ya define \ket y \bra; DeclareRobustCommand los redefine sin error.
\DeclareRobustCommand{\ket}[1]{\left|#1\right\rangle}
\DeclareRobustCommand{\bra}[1]{\left\langle#1\right|}
\DeclareRobustCommand{\braket}[2]{\left\langle#1\middle|#2\right\rangle}
\newcommand{\C}{\mathbb{C}}
\newcommand{\R}{\mathbb{R}}
\DeclareMathOperator{\Tr}{Tr}

% ---- Cajas de apoyo --------------------------------------------------------
% "concepto": definiciones previas que hacen falta para entender el capitulo.
\newtcolorbox{concepto}[1]{breakable,colback=blue!4,colframe=blue!55!black,
  fonttitle=\bfseries,title={Concepto previo: #1}}
% "nota": aclaraciones y advertencias.
\newtcolorbox{nota}[1]{breakable,colback=orange!6,colframe=orange!70!black,
  fonttitle=\bfseries,title={#1}}

% ---- Estilos de listados ----------------------------------------------------
\lstdefinestyle{py}{language=Python,basicstyle=\ttfamily\footnotesize,
  breaklines=true,postbreak=\mbox{\textcolor{gray}{$\hookrightarrow$}\space},
  numbers=left,numberstyle=\tiny\color{gray},frame=single,
  commentstyle=\color{green!40!black}\itshape,
  keywordstyle=\color{blue!70!black}\bfseries,
  stringstyle=\color{red!60!black},showstringspaces=false,tabsize=4,
  columns=fullflexible,keepspaces=true}
\lstdefinestyle{salida}{basicstyle=\ttfamily\footnotesize,breaklines=true,
  frame=leftline,backgroundcolor=\color{gray!8},columns=fullflexible,
  keepspaces=true}

  \begin{document}
  \title{\tituloconcepto}
  \author{\autor}
  \date{\today}
  \maketitle
\fi
 cada capitulo define \tituloconcepto.
\ifdefined\enmain\else
  \documentclass[11pt,a4paper]{article}
  % Preambulo compartido por todos los documentos. Compilar con XeLaTeX/LuaLaTeX
% (o Tectonic). Si se usa Overleaf, elegir el compilador XeLaTeX.
\usepackage[spanish,es-noshorthands]{babel}
\usepackage{amsmath,amssymb}
\usepackage{booktabs,array}
\usepackage[a4paper,margin=2.5cm]{geometry}
\usepackage{xcolor}
\usepackage{listings}
\usepackage[most]{tcolorbox}
\usepackage{quantikz}
\usepackage[hidelinks]{hyperref}

% ---- Autor (editar a gusto) -------------------------------------------------
\newcommand{\autor}{Genesis Silesky}

% ---- Notacion de Dirac y atajos ---------------------------------------------
% quantikz ya define \ket y \bra; DeclareRobustCommand los redefine sin error.
\DeclareRobustCommand{\ket}[1]{\left|#1\right\rangle}
\DeclareRobustCommand{\bra}[1]{\left\langle#1\right|}
\DeclareRobustCommand{\braket}[2]{\left\langle#1\middle|#2\right\rangle}
\newcommand{\C}{\mathbb{C}}
\newcommand{\R}{\mathbb{R}}
\DeclareMathOperator{\Tr}{Tr}

% ---- Cajas de apoyo --------------------------------------------------------
% "concepto": definiciones previas que hacen falta para entender el capitulo.
\newtcolorbox{concepto}[1]{breakable,colback=blue!4,colframe=blue!55!black,
  fonttitle=\bfseries,title={Concepto previo: #1}}
% "nota": aclaraciones y advertencias.
\newtcolorbox{nota}[1]{breakable,colback=orange!6,colframe=orange!70!black,
  fonttitle=\bfseries,title={#1}}

% ---- Estilos de listados ----------------------------------------------------
\lstdefinestyle{py}{language=Python,basicstyle=\ttfamily\footnotesize,
  breaklines=true,postbreak=\mbox{\textcolor{gray}{$\hookrightarrow$}\space},
  numbers=left,numberstyle=\tiny\color{gray},frame=single,
  commentstyle=\color{green!40!black}\itshape,
  keywordstyle=\color{blue!70!black}\bfseries,
  stringstyle=\color{red!60!black},showstringspaces=false,tabsize=4,
  columns=fullflexible,keepspaces=true}
\lstdefinestyle{salida}{basicstyle=\ttfamily\footnotesize,breaklines=true,
  frame=leftline,backgroundcolor=\color{gray!8},columns=fullflexible,
  keepspaces=true}

  \begin{document}
  \title{\tituloconcepto}
  \author{\autor}
  \date{\today}
  \maketitle
\fi


\section{Algoritmo de Grover: búsqueda cuadrática}

\subsection{Objetivo}
Implementar Grover para $N=8$ elementos y comparar el número de consultas con
la búsqueda clásica, dejando claro que la ventaja es \emph{cuadrática}, no
exponencial.

\subsection{Conceptos previos}
\begin{concepto}{oráculo, reflexión y complejidad de consultas}
\begin{itemize}
  \item \textbf{Oráculo:} una caja negra $O$ que reconoce la solución. Aquí
    $O\ket x=-\ket x$ si $x$ es el elemento marcado y $O\ket x=\ket x$ si no.
  \item \textbf{Complejidad de consultas:} se cuenta cuántas veces se usa el
    oráculo, no el tiempo total. Clásico: $O(N)$; Grover: $O(\sqrt N)$.
  \item \textbf{Interferencia} (capítulo~05a): Grover amplifica la amplitud de
    la respuesta correcta.
\end{itemize}
\end{concepto}

\subsection{Fundamento teórico}
Se parte de la superposición uniforme
$\ket s=H^{\otimes n}\ket{0}^{\otimes n}=\tfrac1{\sqrt N}\sum_x\ket x$ con
$N=2^n$. Cada iteración aplica $G=D\,O$, donde $D=2\ket s\bra s-I$
(``difusión'': reflexión respecto al promedio de las amplitudes). Geométricamente
$G$ gira el estado, en el plano generado por $\ket s$ y la solución $\ket w$,
un ángulo $2\theta$ con
\begin{equation}
  \sin\theta=\frac1{\sqrt N}.
\end{equation}
Tras $k$ iteraciones, la probabilidad de medir la solución es
\begin{equation}
  P_k=\sin^2\!\big((2k+1)\theta\big),\qquad
  k_{\text{opt}}\approx\frac\pi4\sqrt N .
\end{equation}
Para $N=8$: $\theta=\arcsin(1/\sqrt8)\approx0.3614$ y
$k=\lfloor\tfrac\pi4\sqrt8\rfloor=2$, de donde
$P_2=\sin^2(5\theta)\approx0.945$.

\paragraph{Contraste.} Clásicamente, buscar un elemento entre $N$ sin
estructura requiere $\sim N/2$ consultas en promedio; Grover usa
$\sim\tfrac\pi4\sqrt N$ (mejora cuadrática, que se sabe óptima para este
problema). Esto es muy distinto de Shor (factorización), cuya ventaja sobre
los mejores algoritmos clásicos conocidos es \emph{exponencial}.

\paragraph{Oráculo en el código.} Para marcar $\ket{101}$ se aplica $X$ a los
qubits donde esa cadena tiene un $0$ (convirtiéndola en $\ket{111}$), luego un
$CCZ$ (hecho como $H$--Toffoli--$H$ en el último qubit), y se deshacen las $X$.
El $CCZ$ solo cambia el signo cuando los tres qubits valen $1$.

\subsection{Código (completamente comentado)}
\lstinputlisting[style=py]{../06_grover_vs_clasico.py}

\subsection{Explicación celda por celda}
\begin{description}
  \item[Celda 2] Construye el oráculo. \texttt{reversed(marcado)} se usa
    porque Qiskit numera el qubit $0$ a la derecha. \texttt{mcx} es una puerta
    X con varios controles.
  \item[Celda 3] \texttt{grover\_operator} arma $G$ a partir del oráculo.
    Se calcula $k_{\text{opt}}$, se aplica $H$ a todos y se repite $G$
    $k$ veces. \texttt{Statevector} da la probabilidad exacta de
    \texttt{"101"}.
  \item[Celda 4] Tabla comparativa de consultas: la ventaja crece con $N$ pero
    solo como raíz.
\end{description}

\subsection{Salida obtenida}
\lstinputlisting[style=salida]{../salidas/06_grover_vs_clasico.txt}
Coincide con la teoría: $P_2\approx0.945$.

\subsection{Interpretación}
Con $N=10^{12}$ Grover necesita unas $7.9\times10^5$ consultas frente a
$5\times10^{11}$: una mejora enorme, pero cuadrática. Además cada consulta
cuántica puede ser mucho más cara que una clásica, y el oráculo debe
construirse como un circuito; por eso la ventaja práctica depende del
problema y del hardware.

\subsection{Ejercicios}
\begin{enumerate}
  \item Cambie \texttt{marcado} a \texttt{"010"} y verifique el resultado.
  \item Use 1 y 3 iteraciones en lugar de 2: ¿qué pasa con $P_k$? (Pista:
    es una función periódica de $k$; ``pasarse'' empeora el resultado.)
  \item Extienda a $n=4$ y calcule el $k_{\text{opt}}$.
\end{enumerate}

\ifdefined\enmain\else
  \end{document}
\fi

